%%FRAG s:codec
%%BEFORE: \subsubsection*{Khai báo mức module (dòng 12--25)}
\subsubsection{Cơ sở toán: Giải mã và mã hóa file PLY}\label{s:codec}
Một gaussian gồm tâm $\mu$, màu $c$, độ đục $\alpha$, kích thước $s$, quay $q$. File lưu giá trị ``thô'':

\begin{center}\small
\begin{tabular}{@{}lll@{}}\toprule
Thuộc tính & Giải mã & Lý do\\midrule
\code{x,y,z} & $\mu=(x,y,z)$ & \
\code{f_dc_0..2} & $c = 0.5 + C_0\,f_{dc}$,\; $C_0=\tfrac{1}{2\sqrt\pi}=0.28209479$ & $C_0=Y_{00}$ (SH bậc 0)\
\code{opacity} & $\alpha = 1/(1+e^{-\mathrm{opacity}})$ & tối ưu logit để $\alpha\in(0,1)$\
\code{scale_i} & $s_i = \exp(\mathrm{scale}_i)$ & tối ưu log để $s>0$\
\code{rot_0..3} & $q\leftarrow q/\|q\|$; nếu $\|q\|=0$ thì $(1,0,0,0)$ & quaternion phải đơn vị\\bottomrule
\end{tabular}
\end{center}

\textbf{Mã hóa ngược:} $f_{dc}=(c-0.5)/C_0$, $\mathrm{scale}=\ln\max(s,10^{-12})$, $\mathrm{opacity}=\ln\frac{a}{1-a}$ với $a=\mathrm{clip}(\alpha,10^{-6},1-10^{-6})$ để tránh $\ln 0$.

\textbf{Point cloud} (không có \code{scale_0}): bán kính mỗi điểm $=\tfrac12 (V/N)^{1/3}$, với $V$ là tích ba cạnh hộp bao. Đó là ước lượng khoảng cách trung bình giữa các điểm. Màu \code{uchar} chia 255.

\textbf{Dung lượng:} $3+3+3+45+1+3+4=62$ float32 $=248$ byte/gaussian (SH bậc 3: 15 hệ số $\times$ 3 kênh $=45$).
%%END

%%FRAG s:layer
%%BEFORE: \subsubsection*{Khai báo mức module (dòng 39--51)}
\subsubsection{Cơ sở toán: Biến đổi layer}\label{s:layer}
Phần này giải thích toán của \code{core/layer.py}, cùng với \code{core/transform.py} (dữ liệu \code{Transform}) và \code{core/math3d.py} (các hàm quay, xem mục \ref{s:quat}).

\paragraph{Ma trận quay từ góc Euler (\code{euler_to_matrix}).}
\[
R_x=\begin{bmatrix}1&0&0\\0&c_x&-s_x\\0&s_x&c_x\end{bmatrix},\;
R_y=\begin{bmatrix}c_y&0&s_y\\0&1&0\\-s_y&0&c_y\end{bmatrix},\;
R_z=\begin{bmatrix}c_z&-s_z&0\\s_z&c_z&0\\0&0&1\end{bmatrix},\;
R=R_zR_yR_x
\]
Với vector cột, $X$ được áp dụng trước, rồi $Y$, rồi $Z$.

\paragraph{Ngược lại (\code{matrix_to_euler}).}
Từ $R=R_zR_yR_x$ có $R_{20}=-\sin y$, nên $y=\operatorname{asin}(-R_{20})$ (kẹp $[-1,1]$ chống lỗi làm tròn), $x=\operatorname{atan2}(R_{21},R_{22})$, $z=\operatorname{atan2}(R_{10},R_{00})$. Khi $|R_{20}|\ge0.999999$ ($y=\pm90^\circ$, \textbf{khóa gimbal}) thì $x,z$ không xác định riêng, đặt $x=0$, $z=\operatorname{atan2}(-R_{01},R_{11})$.

\paragraph{Công thức Rodrigues (\code{axis_angle_matrix}).}
Quay góc $a$ quanh trục đơn vị $(x,y,z)$, $c=\cos a$, $s=\sin a$, $t=1-c$:
\[
R=\begin{bmatrix}tx^2+c & txy-sz & txz+sy\ txy+sz & ty^2+c & tyz-sx\ txz-sy & tyz+sx & tz^2+c\end{bmatrix}
= I + s[k]_\times + t[k]_\times^2 .
\]

\paragraph{Ma trận thế giới của layer (\code{Layer.world_transform / world_matrix}).}
\[
p_w = s\,A\,p_f + t,\quad A=R(\mathrm{euler})\,F,\quad F=\mathrm{diag}(-1,-1,1),\quad
\Sig_w = s^2 A\Sig_f A^{\top},\quad
M=\begin{bmatrix}sA & t\ 0 & 1\end{bmatrix}.
\]
\begin{itemize}
\item $F$ là phép \textbf{xoay $180^\circ$ quanh $Z$} ($\det=+1$, không phải phản xạ). Nó đưa dữ liệu COLMAP (Y hướng xuống) về Y hướng lên.
\item Chỉ $M$ thay đổi (dữ liệu gaussian giữ nguyên), nên kéo/xoay/zoom layer là tức thời. $M$ được gửi lên GPU thành 4 texel (\code{splat_renderer.set_transforms}, chuyển vị để ra 4 cột).
\end{itemize}

\paragraph{Xoay quanh điểm pivot (\code{Transform.rotate_about(Δ, c)}).}
$R'=\Delta R$, $\;t'=\Delta(t-c)+c$ (dời pivot về gốc, xoay, dời lại). Góc Euler mới lấy bằng \code{matrix_to_euler}$(\Delta R)$.

\paragraph{Xuất file (\code{Layer._export_block}).}
Đổi từ thế giới về file: $\mathrm{rot}'=F\cdot\mathrm{rot}$, $t'=F t$. Khi đó:
\begin{itemize}
\item vị trí $p'=(s\,\mathrm{rot}')p+t'$;
\item pháp tuyến $n'=\mathrm{rot}'\,n$ (xoay thuần, không co giãn);
\item scale (log): $\mathrm{scale}_i'=\mathrm{scale}_i+\ln s$, vì $\exp(\mathrm{scale}+\ln s)=s\exp(\mathrm{scale})$;
\item quaternion $q'=q_{\mathrm{rot}}\otimes q$;
\item SH: $\mathrm{sh}'=D(\mathrm{rot}')\,\mathrm{sh}$ (mục \ref{s:shrot}).
\end{itemize}

\paragraph{\code{Layer.world_splats} (dùng cho vật lý, bounds).}
\code{positions = sp.positions @ (s*rot).T + t}; \code{rotations = quat_mul(q_rot, q)}; \code{scales *= s}.
%%END
